\expandafter\ifx\csname CAPETINCLUDEMODE\endcsname\relax
\documentclass[11pt,a4paper]{article}
% preamble.tex — packages et commandes communes a toutes les fiches CAPET SII IE
\input{glyphtounicode}
\pdfgentounicode=1
\usepackage{cmap}
\usepackage[T1]{fontenc}
\usepackage[utf8]{inputenc}
\usepackage[french]{babel}
\usepackage[a4paper, top=1.6cm, bottom=1.6cm, left=1.4cm, right=1.4cm]{geometry}
\setlength{\headheight}{14.5pt}
\usepackage{amsmath, amssymb, amsfonts}
\usepackage{siunitx}
\sisetup{locale=FR, output-decimal-marker={,}, inter-unit-product=\ensuremath{{}\cdot{}}}
\usepackage{array, booktabs, multirow, tabularx}
\usepackage{xcolor}
\usepackage[most, breakable]{tcolorbox}
\usepackage{multicol}
\usepackage{enumitem}
\usepackage{mathtools}
\usepackage{lmodern}
\usepackage{microtype}
\usepackage{fancyhdr}
\usepackage{tikz}
\usetikzlibrary{positioning,arrows.meta}
\usepackage{lastpage}

% ── Couleurs par domaine ──────────────────────────────────────────────
\definecolor{colElec}{RGB}{26, 82, 148}       % Électrotechnique — bleu foncé
\definecolor{colElecLight}{RGB}{210, 225, 245} % fond bleu clair
\definecolor{colAuto}{RGB}{34, 120, 66}        % Automatique — vert
\definecolor{colMeca}{RGB}{185, 90, 20}        % Mécanique — orange
\definecolor{colTherm}{RGB}{160, 40, 40}       % Thermique — rouge
\definecolor{colQT}{RGB}{232, 244, 233}        % fond question type (vert très léger)
\definecolor{colJury}{RGB}{255, 243, 205}      % fond jury (jaune ambre)
\definecolor{colJuryFrame}{RGB}{200, 130, 0}   % bordure jury
\definecolor{colGrayLight}{RGB}{245, 245, 247} % fond gris neutre
\definecolor{colSeparator}{RGB}{180, 180, 180} % gris séparateur

% ── Style tcolorbox de section principale ────────────────────────────
\tcbset{
  ficheSection/.style={
    breakable,
    enhanced,
    colback=#1!7!white,
    colframe=#1,
    colbacktitle=#1,
    coltitle=white,
    fonttitle=\bfseries\normalsize,
    left=4pt, right=4pt, top=3pt, bottom=4pt,
    before skip=6pt, after skip=6pt,
  }
}

% ── Style encadré question type ──────────────────────────────────────
\tcbset{
  questiontype/.style={
    breakable,
    enhanced,
    colback=colQT,
    colframe=colElec!55!black,
    colbacktitle=colElec!55!black,
    coltitle=white,
    fonttitle=\bfseries\small,
    left=4pt, right=4pt, top=3pt, bottom=3pt,
    before skip=5pt, after skip=5pt,
  }
}

% ── Style encadré jury ───────────────────────────────────────────────
\tcbset{
  jury/.style={
    breakable,
    enhanced,
    colback=colJury,
    colframe=colJuryFrame,
    fonttitle=\bfseries\normalsize,
    left=4pt, right=4pt, top=3pt, bottom=3pt,
  }
}

% ── Mise en page ─────────────────────────────────────────────────────
\pagestyle{fancy}
\fancyhf{}
\renewcommand{\headrulewidth}{0.4pt}
\renewcommand{\footrulewidth}{0.3pt}

\setlist[itemize]{leftmargin=1.2em, itemsep=1pt, topsep=2pt, parsep=0pt}
\setlist[enumerate]{leftmargin=1.2em, itemsep=3pt, topsep=2pt, parsep=0pt}

\setlength{\parskip}{2pt}
\setlength{\columnsep}{10pt}

\def\CAPETBEGINFILE{\begin{document}}
\def\CAPETENDFILE{\end{document}}
\else
\def\CAPETBEGINFILE{}
\def\CAPETENDFILE{}
\fi

\fancyhead[L]{\textbf{\textcolor{colElec}{T22}} - CAN/CNA - Numerisation et echantillonnage}
\fancyhead[R]{\small\textit{CAPET SII IE - Session 2026}}
\fancyfoot[C]{\small\thepage\,/\,\pageref{LastPage}}

\CAPETBEGINFILE

% -----------------------------------------------------------------------
%  BLOC EN-TeTE
% -----------------------------------------------------------------------

\begin{tcolorbox}[
  enhanced, breakable,
  colback=colElec,
  colframe=colElec,
  left=6pt, right=6pt, top=5pt, bottom=5pt,
  before skip=0pt, after skip=4pt
]
  {\color{white}\bfseries\large T22 \quad CAN/CNA - Numerisation et echantillonnage}\\[2pt]
  {\color{colElec!30!white}\small
    \textbf{Domaine :} electronique - Chaene d'acquisition numerique \quad
    \textbf{Presence :} \textbf{8/13} sujets (2013-2026) - theme \emph{tres frequent}
  }\\[1pt]
  {\color{colElec!20!white}\footnotesize
    Sous-themes : quantum - nombre de bits - codage binaire/hex - theoreme de Shannon - filtre anti-repliement - chaene capteureconditionneureCAN - CNA en sortie de correcteur
  }
\end{tcolorbox}

% -----------------------------------------------------------------------
%  1. RAPPELS DE COURS
% -----------------------------------------------------------------------

\begin{tcolorbox}[
  ficheSection=colElec,
  title={1. Rappels de cours},
]

% -- 1.1 CAN - Quantum et resolution -----------------------------------
\textbf{1.1 Convertisseur Analogique-Numerique (CAN) - Quantum}

\[
  q = \frac{V_{ref}}{2^n - 1} \approx \frac{V_{ref}}{2^n} \quad (n \geq 8)
  \qquad
  N = \frac{V_E}{q} = \frac{V_E \cdot 2^n}{V_{ref}}
  \qquad
  \text{Erreur max :} \pm\frac{q}{2}
\]

\textbf{Nombre de bits necessaires} pour une precision $\delta$ sur la plage $V_{ref}$ :
\[
  2^n \geq \frac{V_{ref}}{\delta}
  \qquad \Rightarrow \qquad
  n \geq \log_2\!\left(\frac{V_{ref}}{\delta}\right)
\]

\begin{center}
\renewcommand{\arraystretch}{1.25}
\begin{tabular}{cccc}
\toprule
$n$ (bits) & $2^n$ & $q$ pour $V_{ref}=5\,\mathrm{V}$ & $q$ pour $V_{ref}=10\,\mathrm{V}$ \\
\midrule
8  & 256   & $\approx 19{,}6\,\mathrm{mV}$ & $\approx 39{,}2\,\mathrm{mV}$ \\
10 & 1\,024 & $\approx 4{,}9\,\mathrm{mV}$  & $\approx 9{,}8\,\mathrm{mV}$ \\
12 & 4\,096 & $\approx 1{,}2\,\mathrm{mV}$  & $\approx 2{,}44\,\mathrm{mV}$ \\
\bottomrule
\end{tabular}
\end{center}

\medskip\hrule\vspace{4pt}

% -- 1.2 CNA ----------------------------------------------------------
\textbf{1.2 Convertisseur Numerique-Analogique (CNA)}

\[
  V_S = N \cdot q = N \cdot \frac{V_{ref}}{2^n}
\]
Utilise en sortie du correcteur numerique dans une boucle d'asservissement echantillonne.

\medskip\hrule\vspace{4pt}

% -- 1.3 Shannon ------------------------------------------------------
\textbf{1.3 Theoreme de Shannon - echantillonnage}

\[
  f_e \geq 2\,f_{max}
\]
\begin{itemize}
  \item Filtre anti-repliement (passe-bas) avant le CAN : $f_c \leq f_e/2$
  \item Periode d'echantillonnage en asservissement : $T_e \leq T_{systeme}/10$ (regle pratique)
  \item Aliasing (repliement spectral) si $f_e < 2\,f_{max}$ - signal reconstruit corrompu
\end{itemize}

\medskip\hrule\vspace{4pt}

% -- 1.4 Chaene d'acquisition -----------------------------------------
\textbf{1.4 Chaene d'acquisition complete}

\[
  \text{Grandeur physique} \xrightarrow{\text{Capteur}} v_T
  \xrightarrow{\text{Conditionneur+AMP}} v_A \in [0\,;\,V_{ref}]
  \xrightarrow{\text{CAN}} G_N
\]
La tension $v_A$ doit couvrir \textbf{toute} la plage $[0\,;\,V_{ref}]$ pour utiliser au maximum la resolution.

\textbf{Sensibilite globale de la chaene :}
\[
  S = \frac{\Delta G_N}{\Delta x} \quad [\text{LSB/unite physique}]
\]

\medskip\hrule\vspace{4pt}

% -- 1.5 Codage hexadecimal --------------------------------------------
\textbf{1.5 Codage hexadecimal}

Grouper par 4 bits depuis le bit de poids faible (LSB) : ex. $1842_{10} = 0111\,0011\,0010_2 = 732_{16}$.

\end{tcolorbox}

% -----------------------------------------------------------------------
%  2. FORMULAIRE SYNTHeTIQUE
% -----------------------------------------------------------------------

\begin{tcolorbox}[
  ficheSection=colElec,
  title={2. Formulaire synthetique},
]
\small
\begin{multicols}{2}
\renewcommand{\arraystretch}{1.45}
\begin{tabular}{@{}p{4.8cm}ll@{}}
\toprule
\textbf{Formule} & \textbf{Usage} & \textbf{U.} \\
\midrule
\multicolumn{3}{@{}l}{\footnotesize\textit{CAN}} \\
$q = V_{ref}/(2^n-1)$ & Quantum exact & V \\
$q \approx V_{ref}/2^n$ & Quantum approx. & V \\
$N = V_E/q$ & Code numerique & LSB \\
$n \geq \log_2(V_{ref}/\delta)$ & Nbre de bits & - \\
Erreur max $= \pm q/2$ & Quantification & V \\
\midrule
\multicolumn{3}{@{}l}{\footnotesize\textit{CNA}} \\
$V_S = N \cdot V_{ref}/2^n$ & Reconstruction & V \\
\bottomrule
\end{tabular}

\columnbreak

\renewcommand{\arraystretch}{1.45}
\begin{tabular}{@{}p{4.8cm}ll@{}}
\toprule
\textbf{Formule} & \textbf{Usage} & \textbf{U.} \\
\midrule
\multicolumn{3}{@{}l}{\footnotesize\textit{Shannon}} \\
$f_e \geq 2\,f_{max}$ & echantillonnage & Hz \\
$f_c \leq f_e/2$ & Filtre AR & Hz \\
$T_e \leq T_{syst}/10$ & Asservissement & s \\
\midrule
\multicolumn{3}{@{}l}{\footnotesize\textit{Donnees types}} \\
10 bits / 5\,V : $q \approx 4{,}9\,\mathrm{mV}$ & PIC 18F & - \\
12 bits / 10\,V : $q \approx 2{,}44\,\mathrm{mV}$ & Precision & - \\
S7-1500 : PIW8 $\in [0\,;\,27\,648]$ pour 4-20\,mA & API & - \\
\bottomrule
\end{tabular}
\end{multicols}
\normalsize
\end{tcolorbox}

% -----------------------------------------------------------------------
%  3. QUESTIONS TYPES
% -----------------------------------------------------------------------

\begin{tcolorbox}[
  ficheSection=colElec,
  title={3. Questions types - extraites des sujets reels},
]

% -- QT1 -------------------------------------------------------------

% Mini-procedure de resolution type copie concours
\begin{tcolorbox}[questiontype,title={Procedure type copie concours (5 etapes)}]
\begin{enumerate}[leftmargin=1.5em]
  \item \textbf{Qualifier le modele} : hypotheses, domaine de validite, grandeurs d'entree/sortie.
  \item \textbf{Ecrire les lois} : equations de base avec symboles et unites SI.
  \item \textbf{Isoler l'inconnue} : expression litterale, puis application numerique tracee.
  \item \textbf{Controler l'ordre de grandeur} : coherence physique et unite finale.
  \item \textbf{Conclure concours} : phrase courte orientee decision technique.
\end{enumerate}
\end{tcolorbox}

% Visuel de methode : chemin de resolution CAPET
\begin{center}
\begin{tikzpicture}[font=\small, node distance=7mm and 6mm, >=Stealth]
  \node[draw, rounded corners, fill=colElec!8, inner sep=4pt] (n1) {Notions};
  \node[draw, rounded corners, fill=colElec!8, right=of n1, inner sep=4pt] (n2) {Methode};
  \node[draw, rounded corners, fill=colElec!8, right=of n2, inner sep=4pt] (n3) {Application annale};
  \node[draw, rounded corners, fill=colElec!8, right=of n3, inner sep=4pt] (n4) {Verification jury};
  \draw[->, thick, colElec] (n1) -- (n2);
  \draw[->, thick, colElec] (n2) -- (n3);
  \draw[->, thick, colElec] (n3) -- (n4);
\end{tikzpicture}
\end{center}
\begin{tcolorbox}[questiontype,
  title={QT1 - Calcul du quantum et codage hexadecimal \hfill\normalfont\footnotesize E1 2013 (tour laboratoire)}
]
\textbf{enonce.}
Signal de temperature compris entre $-50^\circ$C et $+50^\circ$C, precision $0{,}1^\circ$C, tension de sortie du capteur-conditionneur : $V \in [0\,;\,10]\,\mathrm{V}$. Determiner le nombre de bits du CAN et calculer le mot hexadecimal pour $\theta = -5^\circ$C.

\tcbline

\textbf{Corrige.}

Precision en tension : $\delta = \dfrac{0{,}1}{100} \times 10 = 10\,\mathrm{mV}$.

\[
  2^n \geq \frac{10}{0{,}01} = 1000 \quad \Rightarrow \quad n \geq 10 \; \Rightarrow \; \text{choix : } \mathbf{n = 12\,\text{bits}}
\]

Quantum : $q = 10/(4095) \approx 2{,}44\,\mathrm{mV}$.

Tension pour $\theta = -5^\circ$C : $V_E = \dfrac{10}{100}(-5+50) = 4{,}5\,\mathrm{V}$.

\[
  N = \frac{4{,}5}{2{,}44 \times 10^{-3}} = 1842_{10} \quad = \quad 0111\,0011\,0010_2 \quad = \quad \mathbf{732_{16}}
\]
\end{tcolorbox}

\vspace{4pt}

% -- QT2 -------------------------------------------------------------

\begin{tcolorbox}[questiontype,
  title={QT2 - Sensibilite de la chaene et mot numerique \hfill\normalfont\footnotesize E2 BAC3 2026 (Q16-Q17)}
]
\textbf{enonce.}
Capteur Jumo + conditionneur + CAN 10 bits, $V_{ref} = 3\,\mathrm{V}$, plage $T \in [-5\,;\,90\mathrm{C}]$. La sensibilite globale est $S = 2{,}684\,\mathrm{LSB/C}$. Calculer le mot numerique $G_N$ pour $T_C = 55^\circ$C.

\tcbline

\textbf{Corrige.}

Verification : plage de $95^\circ$C, $V_{ref} = 3\,\mathrm{V}$, $q = 3/1023 \approx 2{,}933\,\mathrm{mV}$, gain conditionneur $c = 3/95 \approx 0{,}0316\,\mathrm{V/C}$ - $S = c/q = 0{,}0316/2{,}933 \times 10^{-3} = 2{,}684\,\mathrm{LSB/C}$. - 
Pour $T_C = 55^\circ$C (par rapport - $T_{min} = -5^\circ$C) :
\[
  G_N = S \times (T_C - T_{min}) = 2{,}684 \times (55 - (-5)) = 2{,}684 \times 60 \approx \mathbf{161\,\mathrm{LSB}}
\]
En hexadecimal : $161_{10} = 1010\,0001_2 = \mathrm{A1_{16}}$.
\end{tcolorbox}

\vspace{4pt}

% -- QT3 -------------------------------------------------------------

\begin{tcolorbox}[questiontype,
  title={QT3 - Theoreme de Shannon - filtre anti-repliement \hfill\normalfont\footnotesize E1 2026 (B17)}
]
\textbf{enonce.}
Rappeler le theoreme de Shannon. Definir le rele du filtre anti-repliement. Deduire sa frequence de coupure $f_c$ si la frequence d'echantillonnage est $f_e$.

\tcbline

\textbf{Corrige.}

\textbf{Theoreme de Shannon :} Pour reconstituer sans aliasing un signal de frequence maximale $f_{max}$, la frequence d'echantillonnage $f_e$ doit verifier $f_e \geq 2\,f_{max}$.

\textbf{Filtre anti-repliement :} filtre passe-bas place \textbf{avant le CAN}, qui elimine toutes les composantes spectrales superieures - $f_e/2$. Sans ce filtre, ces composantes se replient dans la bande utile (aliasing) et corrompent le signal numerique.

\textbf{Frequence de coupure (filtre parfait) :}
\[
  f_c = \frac{f_e}{2}
\]
En pratique, $f_c < f_e/2$ pour tenir compte de la transition finie du filtre reel.
\end{tcolorbox}

\end{tcolorbox}

% -----------------------------------------------------------------------
%  4. ATTENTES DU JURY
% -----------------------------------------------------------------------

\begin{tcolorbox}[
  jury,
  title={4. Attentes du jury - Erreurs recurrentes et points de vigilance},
]

\begin{enumerate}[label=\textbf{\arabic*.}, leftmargin=*, itemsep=5pt]

  \item \textbf{Maetriser la formule du quantum} et la distinguer de la precision physique. Resolution $= q$ en volt ; precision physique $= q / (\text{sensibilite capteur})$ en unite mesuree.

  \item \textbf{Utiliser $2^n - 1$ ou $2^n$} selon le contexte. Pour $n \geq 8$, l'approximation $2^n$ est admise mais doit etre justifiee.

  \item \textbf{Shannon : ne pas se limiter - l'enonce.} Appliquer la consequence immediate : $f_c \leq f_e/2$ pour le filtre anti-repliement, $T_e \leq T_{syst}/10$ pour l'asservissement.

  \item \textbf{Filtre anti-repliement obligatoire.} Le jury attend sa mention des qu'un CAN apparaet dans une boucle de mesure - meme si la question ne le demande pas explicitement.

  \item \textbf{Codage hexadecimal.} Grouper 4 bits depuis le LSB, pas depuis le MSB. Verifier la coherence : $N_{10} \times q$ doit retrouver $V_E$.

  \item \textbf{Confondre CAN et CNA.} Le CAN est en voie de mesure (entree analogique - code numerique) ; le CNA est en voie de commande (sortie correcteur numerique - actionneurs analogiques).

\end{enumerate}

\medskip
\begin{tcolorbox}[
  colback=colJuryFrame!12!colJury,
  colframe=colJuryFrame,
  left=4pt, right=4pt, top=2pt, bottom=2pt,
  fontupper=\small\itshape
]
Rapport jury 2020 : "La majeure partie des candidats ne sait pas determiner la resolution d'un capteur. Les resultats doivent etre encadres, les unites coherentes, la verification de coherence dimensionnelle obligatoire."
\end{tcolorbox}

\end{tcolorbox}

% -----------------------------------------------------------------------
%  PIED - PReSENCE DANS LES SUJETS
% -----------------------------------------------------------------------

\begin{tcolorbox}[
  enhanced, breakable,
  colback=colGrayLight,
  colframe=colSeparator,
  left=4pt, right=4pt, top=3pt, bottom=3pt,
  before skip=4pt, after skip=0pt,
  fonttitle=\small\bfseries,
  title={Presence dans les sujets analyses (2013-2026)},
]
\begin{center}
\footnotesize\renewcommand{\arraystretch}{1.15}
\begin{tabular}{lccccccccccccc}
\toprule
\textbf{Annee} & 13 & 14 & 15 & 16 & 17 & 18 & 19 & 20 & 21 & 22 & 23 & 24 & 25 \\
\midrule
E1 (commune)   & $\checkmark$ & - & - & - & - & - & - & - & - & $\checkmark$ & - & - & $\checkmark$ \\
E2 discipl. IE & - & - & - & - & - & $\checkmark$ & - & $\checkmark$ & $\checkmark$ & $\checkmark$ & - & - & $\checkmark$ \\
E2 appl. comm. & - & - & - & - & - & - & - & - & - & - & - & - & - \\
\bottomrule
\end{tabular}
\end{center}
\smallskip
\footnotesize\textbf{Sous-themes frequents :}
calcul du quantum - choix du nombre de bits - codage hexa - theoreme de Shannon - filtre anti-repliement - chaene capteureCAN complete - sensibilite globale

\medskip
\textbf{Ordres de grandeur (reperes rapides).}
\begin{itemize}[leftmargin=1.3em]
  \item Reseau BT FR : tension phase-neutre typique \SI{230}{\volt}, tension phase-phase typique \SI{400}{\volt}, frequence \SI{50}{\hertz}.
  \item Courants d'etude frequents en sujet : de quelques \si{\ampere} a quelques dizaines de \si{\ampere}.
  \item En conversion d'energie, un rendement attendu reste generalement entre 0{,}8 et 0{,}98 selon la technologie.
\end{itemize}

\medskip
\textbf{Synthese express (a refaire sans support).}
\begin{enumerate}[leftmargin=1.4em]
  \item Poser les hypotheses du modele et les conventions de signe.
  \item Ecrire 2 a 4 relations cles avec unites.
  \item Realiser un calcul numerique complet et verifier l'ordre de grandeur.
  \item Justifier le resultat par un argument technique de niveau CAPET.
\end{enumerate}

\end{tcolorbox}

\CAPETENDFILE









